\section{Conclusão}
\label{sec:conclusao}

Modelos de linguagem de grande porte estão se tornando infraestrutura de trabalho
para a política de poluição do ar, domínio em que uma resposta imprecisa sobre um
padrão vinculante, uma concentração medida ou uma resposta operacional carrega
consequências diretas de saúde pública no Sul Global, onde a carga de mortalidade
por PM\textsubscript{2,5} é mais pesada~\citep{gbd2019risk,requia2024shortterm}.
Medimos viés geográfico e modulado por persona nesse domínio com um benchmark de
25 países, 14 modelos e quatro idiomas (inglês mais o idioma nativo de nove
países), pontuado contra gabarito ancorado em
fontes oficiais primárias --- por código sempre que a resposta é um valor de
registro, e pela média de um painel de juízes de três fornecedores nos demais
casos.

As mitigações que um órgão pode tentar não funcionam, e a mais intuitiva sai pela
culatra. Perguntar no idioma do próprio país \emph{reduz} a acurácia em \nnativa{} pontos
percentuais, de forma mais acentuada onde o idioma é menos representado nos
corpora de treinamento; a persona de gestor local não estreita a lacuna; e o
modelo de 7B com ajuste regional é o mais fraco dos catorze sistemas. A acurácia também
despenca na recuperação de um valor vinculante, a tarefa que um gestor menos pode
se dar ao luxo de ter errada. Uma lacuna de camada Norte/Sul persiste e é
altamente estruturada: ao devolver o padrão nacional vinculante, a chance de
resposta correta do Sul Global é \nort{} da chance do Norte, e a lacuna encolhe
conforme a resposta depende menos de um fato específico do país, desaparecendo na
única tarefa sem valor de registro a acertar. O gradiente monotônico de
desenvolvimento não atinge significância, de modo que relatamos uma diferença
entre grupos que conseguimos demonstrar e uma forma ao longo da faixa de
desenvolvimento que não conseguimos. Quanto ao mecanismo, o desenho dentro dos
países descarta o canal do corpus da língua para uma lacuna medida em inglês,
mas não consegue separar a cobertura do país do desenvolvimento, e não
reivindicamos mais do que isso.

Para as agências ambientais do Sul Global, a mensagem prática é cautelar. Saídas
de LLM sobre padrões vinculantes e dados medidos devem ser tratadas como rascunhos
conferidos contra o diário oficial, e nenhum contorno barato no nível do prompt ou
do modelo substitui essa verificação --- muito menos trocar para o idioma local,
que por esta evidência piora a resposta. Interpretamos o padrão pelo arcabouço da
colonialidade do saber~\citep{mohamed2020decolonial}: a penalidade recai sobre
os idiomas sub-atendidos nos corpora de treinamento~\citep{joshi2020state}, e
recai com mais força sobre o hindi, o mais escasso dos três; se esses idiomas são
também falados onde a carga de saúde da poluição do ar é maior não é algo que
este desenho meça, e a única observação no nível de país de que dispomos (Índia)
aponta na direção contrária (Seção~\ref{sec:discussao}).

Um resultado metodológico generaliza para além do domínio. Reconstruir duas chaves
de gabarito a partir de registros oficiais e tirar a comparação de faixa do juiz,
passando-a para o código, mudou as conclusões substantivas sobre as mesmas
respostas. Quando um efeito nesta literatura se mostra frágil, a chave de gabarito
merece a auditoria antes da amostra, porque uma chave defeituosa acrescenta ruído
\emph{estruturado} que se concentra exatamente onde o dado oficial é mais difícil
de obter.

Protocolo, prompts, registros de gabarito com fontes oficiais, código de análise e
dados de resposta anonimizados são liberados abertamente, com pipeline de
reprodução em um comando e um gate de auditoria de métodos que recomputa cada
número relatado a partir do dado bruto. O estudo também oferece um modelo
reutilizável para auditorias aplicadas de equidade em LLMs --- gabarito de fonte
oficial, adjudicação por código sempre que a resposta é um número conferível, e
auditabilidade de ponta a ponta --- nos domínios de alto risco do setor público em
que esses sistemas já estão sendo implantados.
